\section*{Abstract}
\begin{spacing}{1.4}

Computer-aided diagnosis in cardiology is limited by two data problems: the scarcity of labelled recordings and the class imbalance of infrequent conditions. To address them, this work proposes a parametric generator of synthetic electrocardiogram (ECG) signals based on the superposition of five asymmetric Gaussian functions, one per wave of the P-QRS-T cycle, plus a baseline term. Unlike deep generative models, each of its twenty-two parameters has a direct physiological interpretation, which restricts the output space to morphologies described by explicit amplitudes, positions and widths.

The generator comprises three modules. The first fits the model beat by beat with the L-BFGS-B algorithm, using an analytical gradient and box constraints. On 2,237 normal beats from record 100 of the MIT-BIH Arrhythmia Database, the reconstruction reached a mean Pearson correlation of 0.9969 and a root mean square error (RMSE) of 0.078 standard deviations of the normalized signal, with no systematic residual bias. The second is a physiological rule engine that sets the PR, QRS and QT intervals and the P-wave duration inside published normal ranges, or outside them, to build a control group and a group of subjects with declared deviations; on a cohort of 500 healthy and 500 deviated subjects, the measured intervals matched the label in every case. The third is a rhythm module that generates the inter-beat interval series by spectral synthesis, with the scaling exponent of detrended fluctuation analysis as an input parameter, and couples it to the morphology by scaling only the P and T waves with heart rate. On fifty synthetic 24-hour series, the exponent was recovered with a mean absolute error of 0.014.

The results show that a parametric model with few functions reproduces the P-QRS-T morphology in normal sinus rhythm and yields datasets that are labelled by construction, with the cause of each label traceable. The main limitations are the use of a single record as morphological seed, the representation of the T wave, and the constructive, non-diagnostic nature of the labels.

% PENDING (round 2): add one sentence on the frequency-domain validation (OE5)
% and the evaluation with the pretrained ECGFounder model (OE6).

\vspace{0.5cm}

\noindent\textbf{Keywords:} Synthetic ECG, Gaussian models, Numerical optimization, Detrended fluctuation analysis, Heart rate variability, Data augmentation.

\end{spacing}
